% TOP_PROBLEM: {"id": 30006952, "title": "Leray's Problem for Stationary Navier-Stokes under the Total-Flux Condition in Three Dimensions", "category_id": 9, "category": {"id": 9, "name": "pde", "display_name": "Partial Differential Equations", "description": "PDEs and their applications in physics and geometry.", "slug": "pde", "order_index": 9, "created_at": "2026-07-31T15:26:25.671Z"}, "status": "open"}
% FIRST_POSTED: 2026-09-27
% !TeX program = lualatex
% ATTEMPT_STATUS: unresolved
\documentclass[11pt]{article}
\usepackage[margin=25mm]{geometry}
\usepackage{fontspec}
\setmainfont{Latin Modern Roman}
\usepackage{amsmath,amssymb,mathtools}
\usepackage{unicode-math}
\setmathfont{Latin Modern Math}
\usepackage{xurl,hyperref}
\hypersetup{colorlinks=true,urlcolor=blue}
\setcounter{secnumdepth}{0}
\setlength{\parindent}{0pt}
\setlength{\parskip}{5pt}
\setlength{\emergencystretch}{3em}
\begin{document}
\section*{\texorpdfstring{Leray's Problem for Stationary Navier-Stokes under the Total-Flux Condition in Three Dimensions}{Leray's Problem for Stationary Navier-Stokes under the Total-Flux Condition in Three Dimensions}}\label{30006952}
\subsection{Definitions and mathematical statement}
Let $\Omega\subset{\mathbb{R}}^3$ be bounded and multiply connected, with outward normal $n$, and prescribe velocity trace $a$ on its boundary. The stationary incompressible equations are
\[
 -\nu\Delta u+(u\cdot\nabla)u+\nabla p=0,
 \qquad\nabla\cdot u=0,\qquad u|_{\partial\Omega}=a,
 \quad\nu>0.
\]
In the usual finite-energy formulation, $u\in H^1(\Omega;{\mathbb{R}}^3)$ has trace $a$ and satisfies
\[
 \nu\int_\Omega\nabla u:\nabla\phi
 +\int_\Omega (u\cdot\nabla)u\cdot\phi=0
\]
for smooth compactly supported divergence-free test fields $\phi$. The question asks whether the single compatibility condition $\int_{\partial\Omega}a\cdot n=0$ suffices for weak existence, allowing nonzero flux through individual boundary components and without symmetry.

The record does not quantify boundary regularity or the precise trace space for $a$. Those hypotheses must be taken from the cited existence problem; the equations above do not assert a theorem for arbitrary rough domains or arbitrary boundary distributions.
\par\smallskip\textit{Formulation and concept references:} \href{https://ejde.math.txstate.edu/Volumes/2023/26/foreidooni.pdf}{[S1]}.

\subsection{Short English statement}
Does zero total boundary flux suffice for a stationary three-dimensional Navier--Stokes weak solution, even when separate boundary components carry nonzero flux?
\subsection{Research attempt: a coercive subcase and arbitrary-flux exact solutions}
\textbf{Status and regularity convention.} No existence theorem for unrestricted three-dimensional total-flux data is established here. For the analysis below, $\Omega$ is a bounded connected $C^2$ domain with finitely many boundary components, and the prescribed trace has an $H^1$ solenoidal lift $A$. These are explicit hypotheses of our partial theorem. Smooth traces on smooth domains are sufficient for the explicit examples; no assertion about arbitrary boundary distributions is intended.

\textbf{Source audit.} The introduction of [S1], pages 2--3, describes the inhomogeneous boundary problem on a smooth multiply connected domain and separates the unrestricted three-dimensional question from solved planar and symmetry cases. Its own principal minimax theorem is formulated for homogeneous boundary data, with an additional condition on a convex set. That theorem is not a direct resolution of the stated inhomogeneous problem. The primary paper of Korobkov, Pileckas, and Russo, arXiv:1302.0731, likewise specifies plane and axially symmetric spatial domains. The present calculation does not discard those dimensional or symmetry restrictions from the cited existence results.

\subsubsection{1. Compatibility and the role of a lift}
For an $H^1$ divergence-free velocity with trace $a$, the divergence theorem gives the necessary condition
\[
 \sum_{j=0}^N\Phi_j=0,\qquad
 \Phi_j=\int_{\Gamma_j}a\cdot n\,dS,
                                                               \tag{1}
\]
where $\partial\Omega=\Gamma_0\cup\cdots\cup\Gamma_N$. This is one scalar equation. It does not imply the $N+1$ separate equations $\Phi_j=0$.

Let
\[
 V=\{v\in H^1_0(\Omega;\mathbb R^3):\operatorname{div}v=0\},
 \qquad \|v\|_V=\|\nabla v\|_2.
\]
Poincare's inequality makes this a Hilbert norm. Write $u=A+v$, with $v\in V$, where $A\in H^1$ is divergence-free and has trace $a$. A general construction of a solenoidal extension is a separate linear boundary-value problem. Our nonlinear theorem below takes such an $A$ as input, so it does not smuggle existence of a specially small extension into the argument.

Set
\[
 b(w,z,\phi)=\int_\Omega (w\cdot\nabla)z\cdot\phi.
\]
Hölder's inequality and $H^1\hookrightarrow L^4$ in dimension three make the relevant trilinear expressions continuous. If $w$ is divergence-free and $\phi,z$ vanish on the boundary when needed, integration by parts gives
\[
 b(w,z,\phi)=-b(w,\phi,z).
                                                               \tag{2}
\]
The identity follows first for smooth fields and then by Sobolev approximation. In particular,
\[
 b(v,v,v)=0,\qquad b(A,v,v)=0.                            \tag{3}
\]
The second equality remains valid even when $A\cdot n$ is nonzero: the boundary term contains $|v|^2$, whose trace is zero. Confusing that cancellation with $b(v,A,v)=0$ would remove the central difficulty incorrectly.

The weak equation for $v$ is
\[
 \nu(\nabla v,\nabla\phi)+b(v,v,\phi)+b(A,v,\phi)
 +b(v,A,\phi)
 =-\nu(\nabla A,\nabla\phi)-b(A,A,\phi),\quad\phi\in V.
                                                               \tag{4}
\]
Taking $\phi=v$ formally, or within a finite Galerkin system rigorously, yields
\[
 \nu\|\nabla v\|_2^2
 =-\nu(\nabla A,\nabla v)-b(v,A,v)-b(A,A,v).               \tag{5}
\]
Only the transport of $v$ by $A$ cancels automatically.

\subsubsection{2. A complete small-lift existence theorem}
Let $C_4$ satisfy $\|v\|_4\leq C_4\|\nabla v\|_2$ for $v\in H^1_0(\Omega)$. Suppose
\[
 \kappa:=\nu-C_4\|A\|_4>0.                              \tag{6}
\]
Then a weak solution $u=A+v$ exists with
\[
 \|\nabla v\|_2\leq
 \frac{\nu\|\nabla A\|_2+\|A\|_4^2}{\kappa}.             \tag{7}
\]
In particular, the theorem allows nonzero component fluxes whenever an admissible lift satisfies (6).

Here is the a priori estimate. By (2),
\[
 |b(v,A,v)|=|b(v,v,A)|
 \leq \|A\|_4\|v\|_4\|\nabla v\|_2
 \leq C_4\|A\|_4\|\nabla v\|_2^2,                       \tag{8}
\]
and
\[
 |b(A,A,v)|=|b(A,v,A)|\leq\|A\|_4^2\|\nabla v\|_2.       \tag{9}
\]
The integration by parts in (9) again has zero boundary term because $v$ has zero trace. Substitution in (5) gives $\kappa t^2\leq Lt$, with $t=\|\nabla v\|_2$ and $L=\nu\|\nabla A\|_2+\|A\|_4^2$, proving (7) for every Galerkin solution.

For completeness, choose a $V$-orthonormal basis and let $V_m$ be the span of its first $m$ elements. Define the continuous finite-dimensional residual $F_m(v)$ by pairing the left minus right side of (4) with each basis vector. On the sphere $\|v\|_V=R$,
\[
 F_m(v)\cdot v\geq\kappa R^2-LR>0
 \quad\text{if }R>L/\kappa.
\]
Brouwer's fixed-point theorem implies a zero in this ball. One elementary formulation is: if $F_m$ had no zero there, the map $v\mapsto-RF_m(v)/|F_m(v)|$ would map the ball into itself and have a fixed point on the sphere. At that point $F_m(v)\cdot v<0$, a contradiction. Thus solutions $v_m$ exist with the uniform bound (7).

Pass to a subsequence with
\[
 v_m\rightharpoonup v\text{ in }H^1_0,
 \qquad v_m\longrightarrow v\text{ in }L^4.
                                                               \tag{10}
\]
The strong convergence follows from the compact Sobolev embedding on a bounded regular three-dimensional domain, since $4<6$. The space $V$ is weakly closed. For a fixed basis test field, integrate convection by parts to put its derivative on the test field. Then
\[
 \|(A+v_m)\otimes(A+v_m)-(A+v)\otimes(A+v)\|_2\to0,
\]
by strong $L^4$ convergence. The viscous term passes to the limit weakly. Thus (4) holds first for the dense union of $V_m$ and then for every $\phi\in V$ by continuity. The trace and incompressibility conditions are preserved. This proves the weak velocity formulation required in the question under (6). The usual distributional pressure reconstruction is not needed for this velocity conclusion.

This is an existence argument, not a uniqueness argument. Nothing in (7) says two solutions for the same large data must coincide. It also does not supply a bound when $\kappa\leq0$; the displayed proof genuinely stops there.

\subsubsection{3. Why total flux does not make this proof uniformly coercive}
It is tempting to say that one can always choose an extension with arbitrarily small $L^4$ norm. Nonzero component flux provides a direct obstruction to precisely that claim.

Fix a component $\Gamma_i$ and choose $\chi_i\in C^\infty(\overline\Omega)$ which equals one near $\Gamma_i$ and zero near every other component. Such a function exists because there are finitely many disjoint compact boundary components. For every solenoidal $H^1$ lift $A$, integration by parts gives
\[
 \int_\Omega A\cdot\nabla\chi_i=\Phi_i.
\]
Consequently,
\[
 \|A\|_4\geq
 \frac{|\Phi_i|}{\|\nabla\chi_i\|_{4/3}}.                 \tag{11}
\]
The right side is fixed by the boundary data and domain, independently of how the extension is chosen. Large opposite fluxes can satisfy (1) while making this lower bound arbitrarily large. Thus total cancellation of boundary fluxes does not force (6).

This does not prove that every possible Leray--Hopf extension inequality fails, nor that large-flux solutions fail to exist. It proves the more specific statement needed to audit this attempt: the sufficient small-$L^4$-lift hypothesis cannot be manufactured for all total-flux-compatible data by an arbitrary choice of lift. A different a priori estimate could succeed even when (6) fails.

There is a related sign check. The quantity $b(v,A,v)$ is the integral of the quadratic form $v^T(\nabla A)v$. Incompressibility gives $\operatorname{tr}\nabla A=0$; a trace-zero symmetric part need not have a favorable sign. An energy argument that replaces this quadratic term by a nonpositive quantity solely from $\operatorname{div}A=0$ is unjustified.

\subsubsection{4. An exact family with arbitrarily large component flux}
Let
\[
 \Omega=\{x\in\mathbb R^3:r_0<|x|<r_1\},\qquad
 u_c(x)=c\frac{x}{|x|^3},\qquad
 p_c(x)=-\frac{c^2}{2|x|^4},                              \tag{12}
\]
where $0<r_0<r_1$ and $c\in\mathbb R$ is arbitrary. These fields are smooth up to both boundary spheres.

Since $u_c=\nabla(-c/|x|)$ and $|x|^{-1}$ is harmonic away from the origin,
\[
 \operatorname{div}u_c=0,\qquad \Delta u_c=0,
 \qquad \operatorname{curl}u_c=0.
\]
For a curl-free field, $(u\cdot\nabla)u=\nabla(|u|^2/2)$. Here $|u_c|^2=c^2|x|^{-4}$, so the convection term is exactly canceled by $\nabla p_c$. Thus (12) solves the stationary equations for every viscosity $\nu>0$, with the trace prescribed by $u_c$ itself.

The outer boundary has outward normal $x/r_1$ and flux $4\pi c$. The outward normal of the domain on its inner boundary is $-x/r_0$, giving flux $-4\pi c$. The total flux is zero, and each component flux can have arbitrarily large absolute value. Reversing $c$ reverses the fluxes while leaving the pressure unchanged.

The strain also makes the sign issue concrete:
\[
 \nabla u_c=\frac{c}{|x|^3}
 \left(I-3\frac{x\otimes x}{|x|^2}\right).
                                                               \tag{13}
\]
For $c>0$ its radial eigenvalue is $-2c/|x|^3$ and its two tangential eigenvalues are $c/|x|^3$. It is trace-free and indefinite. Large nonzero flux is therefore compatible with an exact solution while defeating a naive sign or smallness argument. This family rules out treating nonzero component flux itself as an obstruction to existence.

\subsubsection{5. A continuation route and its first missing estimate}
One can attempt to deform boundary data from $0$ to $a$ by $ta$, $0\leq t\leq1$, and use the lift $tA$. The estimate just proved is uniform only while
\[
 \nu-t C_4\|A\|_4>0.
\]
It establishes existence near $t=0$ but supplies no bound on the full interval if the coefficient crosses zero. A continuation or degree argument needs an a priori bound for every solution of its homotopy, not merely existence of one small-parameter branch. Neither the total-flux identity nor the shell example provides such a bound for a general domain.

The finite diagnostic checks (12)--(13) symbolically: divergence, Laplacian, the convective-pressure cancellation, and the radial/tangential strain eigenvalues. The Galerkin compactness proof and the flux obstruction (11) are analytic and do not depend on that computation.

The completed partial results are weak existence under a quantified lift condition, a lower bound preventing that condition from becoming automatic at arbitrary flux, and an exact arbitrary-flux family. The first unresolved step is obtaining a global a priori estimate, or another existence mechanism, for general three-dimensional data satisfying only total flux compatibility. No such step is claimed here.

\subsection{Sources}
\begin{itemize}
\item[S1]Fereidooni, Moameni, and Grewal, Existence of solutions to steady Navier-Stokes equations via a minimax approach, Introduction (2023).
\url{https://ejde.math.txstate.edu/Volumes/2023/26/foreidooni.pdf}.
\textit{Formulation source.}
\item[S2] M. V. Korobkov, K. Pileckas, and R. Russo, Solution of Leray's problem for stationary Navier--Stokes equations in plane and axially symmetric spatial domains, arXiv:1302.0731v1 (2013). Abstract checked for dimension and symmetry scope.
\url{https://arxiv.org/abs/1302.0731v1}.
\end{itemize}
\end{document}
